\chapter[Sermon 35. The Faithful Are Sealed to Salvation, Which They Obtain by the Grace of God in Christ Jesus.]{The Faithful Are Sealed to Salvation, Which They Obtain by the Grace of God in Christ Jesus.}
\label{sermon:035}
\markright{Sermon 35. The Faithful Are Sealed to Salvation, Which They ...}

And I saw another angel ascend from the rising of the sun, who had the seal of the living God, and he cried with a loud voice to the four angels (to whom power was given to hurt the earth and the sea), saying, “Hurt not the earth, neither the sea, neither the trees, until we have sealed the servants of our God in their foreheads.” And I heard the number of those who were sealed, and there were sealed one hundred forty-four thousand of all the tribes of the children of Israel. Of the tribe of Judah were sealed twelve thousand; of the tribe of Reuben were sealed twelve thousand; of the tribe of Gad were sealed twelve thousand; of the tribe of Asher were sealed twelve thousand; of the tribe of Naphtali were sealed twelve thousand; of the tribe of Manasseh were sealed twelve thousand; of the tribe of Simeon were sealed twelve thousand; of the tribe of Levi were sealed twelve thousand; of the tribe of Issachar were sealed twelve thousand; of the tribe of Zebulun were sealed twelve thousand; of the tribe of Joseph were sealed twelve thousand; of the tribe of Benjamin were sealed twelve thousand.

After this I beheld, and lo, a great multitude (which no man could number) of all nations and people, and tongues, stood before the throne and before the Lamb, clothed in long white garments and holding palm branches in their hands, and cried with a loud voice saying: “Salvation be ascribed to him that sits upon the throne of our God, and unto the Lamb.” And all the angels stood around the throne and of the elders and of the four living creatures, and fell down before the throne on their faces, and worshipped God, saying: “Amen. Blessing, and glory, wisdom and thanksgiving, and honour and power, and might, be unto our God for ever and ever. Amen.”

We have heard, brethren, concerning the opening of the sixth seal, that the sun was made black, the moon bloody, and the stars fell from heaven to the earth, and the rest that we have rehearsed. By all of these things was signified the corruption of doctrine. A sorrowful and fearful matter was shadowed with most sorrowful and most terrible parables. We have heard how there followed in the world a great turmoil of things, and with it, grievous despair; and that the winds also were restrained, that they should not blow. But we have experienced how great a grief it is, yea, and destruction also, to lack the air or wind, so much that without breathing and cooling, men must needs wither and be weakened and choked. But with so great an evil are those vexed who are destitute of the preaching of God’s word.

Some person here might say that the entire world will perish in heresies, in the Alcoran, in Papistry, and other corruptions. In what state, then, were our forefathers? Do you think they are all damned? John prevents these things, and with a vision wholly Evangelical—that is to say, with a most profitable consolation—shows that God has an innumerable multitude of them, who even in the midst of those Antichristian times or difficulties are made safe: and that by the mere grace of God, through the intercession of Jesus Christ, of whom alone is salvation: that is to say, whom alone those who are saved may thank for their salvation.

We now have to respond to those of opposing views, who constantly object, arguing that if our ministers are condemned, it would be wicked to condemn all, and therefore they must be saved. They claim they have not heard of our new doctrine, but have maintained the old, and therefore will be saved by adhering to the old ways. To this, we reply that our ministers were saved, we gladly grant and believe it as well; but we add that it is by the free grace of God, as we shall understand more clearly presently, and not by popish superstition. Nor will you be saved by that means; but you must also be saved by Christ, if you desire to be saved.

Rather, considering the singular goodness of God today, the gospel is preached, and is preached even to you, to which you show yourself a rebel, declaring yourself not to be among God’s children, who hear the word of God with joy and keep it. You will have no excuse or pretense to cover your sin. If your ancestors had had the same opportunity that you neglect, good God, how far they would have run before you! Therefore, you knowingly and willingly speak against God and wilfully cast yourself into destruction. Die, then, through your own fault.

Neither does this passage testify only that many will be saved by God’s grace from corruption and in the true faith, even when, in human judgment, there appear to be none or very few faithful; and even when very few or none are saved, because of the exceedingly great corruption of every time? We have also heard and read in the third book of Kings, chapter 19, that Elijah, complaining most grievously of the scarcity of the faithful, understood that God had reserved seven thousand men who had not bowed their knees before Baal. Therefore, the Lord always has his chosen, who, by grace through Christ, are saved in the midst of destruction and perdition.

And the Author of this salvation and preservation is first declared to be an angel ascending from the rising of the sun: namely, the Lord Christ, that sun of righteousness, rising up in those most intense Antichristian darkness, to those who seek God, and driving away the darkness for them. For Christ is the true light of all times, illuminating all those who are illuminated. He gives his people also preachers, who by the word may defend God’s people, that they be not destroyed with that common destruction.

\index{Jerome, Saint}For it is diligently expressed that this angel had a seal, and not a seal only, but the seal of God, and even of the living God. For Christ, who is the image of God unseen, that is to say, the print or express image of his substance, in whom we know, as he himself says to Philip, the Father, has a seal, which is an instrument with which we seal such things as we will have sealed, saved, and confirmed, and discerned from that which is counterfeit, and kept safe against deception. But the Lord has no such seal as we have in this world; but so by a figure is called the spirit of God, with whom he inspires his faithful, by whom he gives also a lively faith, by the word of the living and eternal God. This seal therefore is the seal of the living God, the spirit of life, and lively faith. The Apostle Jerome, speaking: “We also trust in Christ, after the word of truth heard, the gospel of your salvation, wherein after the believed, they were sealed with the holy spirit of promise, \&c.” These things are not divided. For faith is not without the word, nor both these without the holy ghost in the faithful. For Christ works with me by a lawful ministry, by whom he inspires certain that may teach and admonish men, unto whom he gives his faith and spirit, sealing their minds. Christ therefore doth prohibit the ministers of Satan, that they in restraining and letting the free preaching of God’s word, should not proceed to hurt men, before the minds of the chosen be sealed; that is to say, teaches how soever the truth is restrained, and the preaching of the Gospel obscured, yet that the minds of many shall so be furnished with God’s word, and with godly inspiration, which may so live, and be of such efficacy in them, that seducing can either have no place in them, or if it have any at all, can not abide or persevere to the end.

\index{Antichrist}\index{Ezekiel (Book of)}There are also two other places in Scripture testifying that seals were given to people, with which they were marked and delivered from present evil; and they do not contradict our seal of the living God. In Exodus 12, the doorposts of the Israelites were sprinkled with the blood of the lamb. The sign itself would have availed nothing, unless the virtue of God, instituting and consecrating the sign with his word, had turned away the destroying angel; nor has the sign lacked faith when used by God’s holy people. For the godly do not receive God’s ordinances without faith. Therefore, the same power of Christ preserved the Israelites from destruction, which now keeps the faithful from the infection of Antichrist. In Ezekiel 9, one seals the foreheads of the faithful, having the stylus of a scribe and priest. Truly, Christ has always defended his people. And he seals by impressing or writing this mark or letter Tau. That mark signifies, that is to say, the Law, or direction, or Rule. For whoever has the law of God, the word of God, and even the rule of faith engraved in his heart, is safe and secure from all evil. The ancients in old times called the rule of faith and direction, the very articles of the Christian faith—“I believe in God,” and so on. You see, therefore, how all these signs in truth come to one point: those whom the Spirit of God has inspired and illumined with faith by the word are safe and secure from evil. This much concerning the seal.

Now let us also consider who they are that are sealed. We read in Ezekiel, pass through the city of Jerusalem and mark Tau in the foreheads of those mourning and lamenting for all the abominations done in the midst thereof. And here it is said, till we seal the servants of our God. Therefore, the servants of God, and those who are sorry for abominable wickedness, are sealed. The contemners of God—hogs and dogs—are neglected.

Furthermore, it is shown in what part of the body they are sealed. In times past, the blood of the lamb was anointed on the doorposts. In Ezekiel, Tau is marked on their foreheads. Here also, the seal of the living God is imprinted on the foreheads of the faithful. And the forehead represents the mind, the chiefest and most excellent part in man. Spirit and faith are put into the minds of the faithful. Nevertheless, the mark is aptly fixed to the forehead, not to the back, or shoulders. For those enlightened by the word and spirit, and who have faith, do confess the same, and dissemble nothing; and much less are they ashamed, but desire that their glory, which is their faith, might be known of all men. We display most notable things written on our foreheads: that is, most manifest things, whereof we are not ashamed.

If we now apply these things to events that occurred in the past and continue to occur today, they will shed great light upon them. There were found good men, faithful and fearing God, mourning and serving God. And there are still people today, amidst Islam and Roman Catholicism, who explicitly condemn and have condemned this kind of life, openly confessing that it is not the true way of life, that there is no more wicked group of men than their priests, and that they would not entrust themselves or their salvation to them, but rather would consecrate themselves wholly to God’s mercy. And others, who have spent a significant portion of their lives with good zeal, though perhaps not according to knowledge, in those trifles and superstitions, when they reach the end of their lives, despise them altogether; and freely professing the truth, they condemn those trifles and commit themselves wholly to the Christian faith, esteeming nothing more excellent or surer than the rule of faith, which they also desire to hear recited to them as a true confession, and die in the same. All these have the mercy of God sealed with the seal of the living God, and delivers from all taint of the Antichrist and Satan, from corruption and destruction, through Jesus Christ our Lord.

\index{Judgment, Last}But lest we should gather in every age only here one and there one, the Lord himself now makes here a great reckoning, first of the Jews by every tribe he gathers twelve thousand, and after by multiplication, one hundred forty-four thousand; and of the Gentiles a multitude innumerable. Wherefore in every time and age innumerable obtain salvation: however much error, seduction, and destruction reign and rule in the world. These things do highly commend God’s mercy and comfort us exceedingly. And where certain gather hereof that there shall be yet in this world before the judgment a Saturnalian or golden age, wherein these things should be fulfilled, and that all men should come to the kingdom of God, it alludes too much to the gross error of the Millenniums, which is already expelled out of the church of God. These things were fulfilled in old time, and are at this day, and shall be fulfilled likewise, so long as the world shall endure. The kingdom of Satan and of Antichrist shall continue always to the last judgment, and shall still impugn the kingdom of Christ, and seem even to oppress the same; much less ought they to promise us so great security. When the Son of Man shall come, says the Son of Man himself in the Gospel, think you shall he find any faith upon earth? And again: it shall be as in the days of Noah and Lot, the words of the gospel are known, as also those of the blessed Apostles Peter and Paul, 2 Peter 3:1, 1 Thessalonians 4.

\index{Augustine, Saint}\index{Irenaeus}\index{Tertullian}\index{Papias}Those who do not like this explanation of ours briefly argue that the prophecies of the prophets concerning the restoration of Israel have not yet been fulfilled, and that, according to the truth of the eternal God, they must be fulfilled. They suppose, and indeed contend, that a certain or predetermined time must remain, in which all these things may be accomplished. To this I answer plainly, that we shall err shamefully with Papias, Augustine, Irenaeus, Tertullian, and Lactantius, and with those who are called Millenniums, unless we judge rightly here. I believe, therefore, that the same restoration of which the prophets speak must be divided into three times: the first to be historical, extending from King Cyrus to great Pompey, and which Ezra, Nehemiah, and the Author of the book of the Maccabees describe and teach to be fulfilled. The second to begin with the coming of our Savior and proceed until Antichrist and his destruction, which the Apostles and Evangelists have most diligently described, and in which they testify many things to be accomplished. And that the third time should begin from the restored gospel and the last judgment, and continue forevermore: which restitution truly seems to be of all others most perfect and complete, in which God will give to humankind fully what things ever he has promised by the mouths of the prophets and Apostles. Saint Peter has manifestly made mention of this in Acts 3:21, saying: “It behooves Christ to take heaven, until the time of the restoration of all things, which God has spoken by the mouth of all his saints from the time of the prophets.” And the Lord himself in the gospel, speaking of the last judgment, said: “Lift up your heads, because your redemption draws near.”

Perhaps we may divide this matter in this way for greater clarity: the restoration of Israel, or of all the faithful, is certainly either physical or spiritual. The physical restoration may be called historical, and was accomplished by Cyrus, Zerubbabel, Joshua, Ezra, Nehemiah, and the Maccabees. The spiritual restoration is being fulfilled, or will yet be accomplished, by the coming of our wholesome Messiah, our Lord Jesus Christ. And the coming of the Lord is of two kinds: the first, indeed, is in the flesh, in which we believe many things, the Apostles bearing witness to their fulfillment in Christ. In the latter, he shall come again from Heaven for judgment. In that coming, he will most fully accomplish those things that we see as yet unfulfilled. And without doubt, all our hope is directed toward this coming and is comforted by it. Those things spoken by the Apostle in Romans 11 concerning the conversion of the Jews are being fulfilled partly, and partly are being fulfilled daily, and will continue to be fulfilled.

Now we return to the abundance of those who shall be saved and are already saved from the midst of the kingdom of Antichrist, to be declared. John divides the universality of humankind into Jews and Gentiles. Of the Jews, one hundred and forty-four thousand are accounted. And according to our judgment, of a thousand Jews, scarcely one or two seem to be saved; but where, by the testimony of our Savior himself, so great a number is saved, there remains certainly, of this number, an infinite multitude of this stubborn people to be gathered, who shall be saved. And they are not saved by the Law, or by circumcision, or by their damnable obstinacy, but by the grace of God in Christ their Messiah, the only Redeemer, revealed to them mercifully, and received faithfully by them. For if the thief on the cross might be saved, now leaving his life, what shall prevent innumerable Jews from being saved by the same means? Nevertheless, I will determine no measure here. Neither will I by this means frustrate the ministry of the word and Sacraments. However, I know the things here spoken to be true; the measure or manner is known to God, and with him nothing is impossible. And to this serves the Apostle’s doctrine, in 11 to the Romans.

\index{Repentance}This doctrine will make men neglect their own salvation; for already there are those who say, “If the end is well, then all is well.” As though they should say, “Live however you please in this world, immersed in pleasures and bloodshed, and devoted to gluttony, but believe only at the very end of your life, and you shall be saved.” I am not ignorant that there are many unclean pigs and filthy swine, abusing the true doctrine and comfort of the Gospel; but will the abuse of ungodly men take away the truth from us? The children of God, who know that there is no other propitiation or satisfaction for sins but the offering of Christ, therefore do not cease to renew their lives daily by repentance.

Thus, although the faithful doubt nothing but that countless numbers are converted and saved by the Lord at the end of their lives, they do not abuse God’s mercy to license the flesh, but are afraid. For there are other passages that restrain them in order and duty. For the Lord says, “You have been forgiven much, go and sin no more, lest something worse befall you.” Also, “Let us do good while we have time; the time will come when we cannot work.” The parable of the ten virgins declares the same to us. Moreover, if the righteous scarcely be saved, where will the sinner and wicked appear? Furthermore, do not tempt the Lord your God. And countless others of like sort. And when the saints have spent their entire lives behaving blamelessly in God’s righteousness, yet in the last time of their lives they do not trust in that, but in God’s mere mercy through Christ. They always remember how gravely he was rebuked in the Gospel, namely, he who envied the good fortune of him who labored with him in the vineyard, for that he had received so much wages, coming into the vineyard about the last hour of the day, as he had received that who had labored all day long, and again, the thrifty son, for that he was sorry that his wasteful and prodigal brother was received again by his father, and a feast also made for him, and for him who was always obedient and took pains continually, no such thing was prepared.

But the Gentiles he does not arrange into any certain number, but says how he saw a great multitude, which no man could tell; no more than they could the stars, sand, herbs, or grass how many they were in number. He signifies therefore, that in all the world, at all times, innumerable people are saved by Christ. Nevertheless, lest any man should think that it would avail or hinder him to salvation to be born of this or that nation, tribe, or tongue, John adds immediately, “of all tribes, peoples, and tongues to be ordained to salvation indifferently.” Therefore, this difference hinders salvation nothing; but people are found in India, Ethiopia, Barbaria, and in the furthest part of Libya, in Scythia, Tartaria, and in the uttermost ends of the world, who are saved by the grace of Christ.

And because it is difficult to reason about things to come, John here speaks most expressly of those who are not to be saved, but have already achieved salvation and are in Heaven, so that we should not doubt of their salvation. He also depicts the manner of salvation and everlasting blessedness. This treatise refutes those who suppose that souls sleep, nor have the enjoyment of Godhead before the judgment, nor are yet in Heaven. First, he says how they stand before the throne and in the sight of the Lamb. For the first blessedness in the blessed life is to see God as he is and to enjoy his glory, to be with Christ in glory. John 17:1. John 3. White stoles are the garments of triumphant and clean people. As will be declared more at large hereafter, and has been noted once or twice before. It signifies that the blessed souls are adorned with light, and so forth. And the palm also is a token of victory. Pliny treats much of the palm in chapter 4 of book 13. All writers say that the palm was the most ancient badge of a conqueror. And why this tree was chiefly chosen for this use, Aulus Gellius shows the cause in chapter 6 of book 3 of Noct. Att. Writing that in a palm tree there is a certain peculiar thing which agrees with the nature of stout and noble people. For if you lay great weights upon the wood thereof, the palm gives not place downward, but rises up against the weight and bears upward. And for this he cites the authority of Aristotle and Plutarch, to whom you may add also Pliny, book 16, chapter 24.

\index{Erasmus of Rotterdam}All of these things are accompanied by an exceedingly great rejoicing, whereby not only do they give God thanks and praise his mercy, but also openly demonstrate and testify whom they may thank for their salvation. And they say, “Salvation to him,” and so forth. For that is the more accurate rendering, as Erasmus has also noted. They signify that God is blessed not only in himself, but also for having communicated this salvation to them and saved them. Of the throne or seat of God, we spoke before in chapter 4. God the Father himself sits on the throne. It is therefore a phrase of speech, which has this meaning: we owe this our salvation and blessedness to our God, who sits in his throne.

Again they communicate this salvation to the Lamb, that is to Christ. For God by his grace through Christ saves the believers. And since Christ is called the Lamb, the whole mystery of the incarnation and redemption is remembered in the word Lamb; that being indeed reconciled to God by the blood of the host, we are now the heirs of God and the sons of God. \&c. Therefore, the saints in heaven, and our fathers already saved and dwelling in heaven, do testify, and in testifying teach, that they are justified and saved not by Islam, or Roman Catholicism, or any other observances, but by the mere grace of God in Christ.

\index{Papacy}Hereby are refuted two opinions, deeply harmful to the whole world. The first claims that Papists are saved because of their simplicity and strict discipline. For they say that because Papists are ignorant of better things, and perform their works with good intentions, they are saved by those same works. This is most vain and most ungodly. They add that unless we judge this way, surely not one Papist would be saved. I say plainly, no one is saved by Papistry, no more than by Mahometry. For it is called the way of perdition even by Saint Peter himself.

However, I do not think that no one among the Papists is saved. But I believe that countless numbers, as I said before, have at length seen the corruption of Papistry through God’s illumination, have forsaken Papistry, and embraced the sincere Gospel, and so are saved by Christ alone.

The latter assumes that every person, regardless of their religion, will be saved. Against this, the faithful here proclaim: those who are saved are saved by the grace of God through Christ. Therefore, no other religion offers salvation. There is no other name given to people, in which they must be saved, but that of Christ Jesus. No other way is open into heaven, nor any other door: whoever claims otherwise is called by the truth a thief and a murderer. Indeed, they utterly abolish Christ and the entire scripture, those who contend that every person is saved by their own religion. I cannot imagine any doctrine so harmful. Therefore, let us maintain what all the faithful in heaven have taught us: that salvation is from God through Christ.

All the angels in heaven confirm these things, lest anything should lack that belongs to a sure and certain testimony; and also teach us by their example what we should do. They sing together “Amen,” thereby also testifying that salvation is by grace alone through Christ. Again, they fall down and worship God. But how much more ought we, mortal men, by worshipping to attribute this honor to him? And by singing a hymn, they exhibit to us a form of serving God, finally of judging rightly of God, that we attribute nothing to any creature to the reproach of the creator, which belongs to God alone, but ascribe all things to God wholly. The words of this hymn are explained in chapters four and five, so I need not linger on them here. They use “blessing” for “praise”; the rest of the words are plain.

And now let us learn, being taught by so many testimonies and examples of all the saints, forsaking all vain and wicked opinions, to give all glory to God through Christ: to whom be praise and thanks giving. Amen.
